% Общие параметры оформления лабораторной работы.
\usepackage[T2A]{fontenc}
\usepackage[utf8]{inputenc}
\usepackage[russian]{babel}
\usepackage{amssymb,amsmath,mathtools}
\usepackage{graphicx}
\usepackage[table]{xcolor}
\usepackage{booktabs,tabularx,array}
\usepackage{enumitem}
\usepackage{needspace}
\usepackage{fancyhdr}
\usepackage{titlesec}
\usepackage{tcolorbox}
\tcbuselibrary{breakable,skins}
\usepackage[hidelinks]{hyperref}

% Размер наборной полосы и интервалы, согласованные с ЛР1 и ЛР2.
\setlength{\textwidth}{16cm}
\setlength{\textheight}{24cm}
\setlength{\oddsidemargin}{0mm}
\setlength{\evensidemargin}{0mm}
\setlength{\topmargin}{-22mm}
\setlength{\headheight}{15pt}
\setlength{\headsep}{7mm}
\setlength{\footskip}{11mm}
\setlength{\parindent}{0pt}
\setlength{\parskip}{2pt}
\tolerance=10000
\emergencystretch=1.2em
\clubpenalty=10000
\widowpenalty=10000
\displaywidowpenalty=10000
\setlength{\abovedisplayskip}{8pt plus 2pt minus 2pt}
\setlength{\belowdisplayskip}{8pt plus 2pt minus 2pt}
\setlength{\abovedisplayshortskip}{5pt}
\setlength{\belowdisplayshortskip}{6pt}
\setlist[enumerate]{leftmargin=2.7em,topsep=5pt,itemsep=4pt,parsep=0pt}
\setlist[itemize]{leftmargin=2.7em,topsep=4pt,itemsep=3pt,parsep=0pt}
\renewcommand{\arraystretch}{1.12}
\setlength{\tabcolsep}{5pt}
\titleformat{\section}{\normalfont\bfseries\large}{\thesection.}{.5em}{}
\titlespacing*{\section}{0pt}{14pt plus 2pt minus 2pt}{8pt}
\titleformat{\subsection}{\normalfont\bfseries\normalsize}{\thesubsection.}{.5em}{}
\titlespacing*{\subsection}{0pt}{12pt plus 2pt minus 2pt}{6pt}
\titleformat{name=\subsection,numberless}{\normalfont\bfseries\normalsize}{}{0pt}{}

% Колонтитулы и первая страница.
\newcommand{\labtitle}{Две стадии: от скрытого времени к прогнозу}
\newcommand{\labtag}{Лабораторная работа~№~3-4}
\newcommand{\itmoidentity}{\includegraphics[width=22mm]{../branding/itmo_logo.pdf}}
\pagestyle{fancy}
\fancyhf{}
\fancyhead[L]{\includegraphics[width=13mm]{../branding/itmo_logo.pdf}}
\fancyhead[R]{\small Теория вероятностей}
\fancyfoot[C]{\small\thepage}
\renewcommand{\headrulewidth}{0pt}
\renewcommand{\footrulewidth}{0pt}
\fancypagestyle{firstpage}{\fancyhf{}\renewcommand{\headrulewidth}{0pt}\fancyfoot[C]{\small\thepage}}
\newcommand{\dochead}[1]{%
  \noindent\begin{minipage}[c]{.77\textwidth}
   {\small Теория вероятностей\par}\vspace{4pt}
   {\bfseries\large \labtag\par}\vspace{2pt}
   {\bfseries\large \labtitle\par}
  \end{minipage}\hfill
  \begin{minipage}[c]{.20\textwidth}\raggedleft\itmoidentity\end{minipage}\par
  \vspace{10pt}
}

% Краткие примечания студенческой версии.
\newenvironment{studentnote}{\par\medskip\begingroup\small\noindent}{\par\endgroup\medskip}
\newcommand{\stageoutcome}[1]{\par\medskip{\small\textbf{Результат этапа.} #1}\par\medskip}

% Смысловые пометки для преподавателя.
\definecolor{explainbg}{RGB}{239,246,250}
\definecolor{solvebg}{RGB}{238,247,240}
\definecolor{commentbg}{RGB}{253,248,235}
\newtcolorbox{teacherbox}[1][]{enhanced,breakable,sharp corners,boxrule=0pt,frame hidden,colback=explainbg,left=8pt,right=8pt,top=6pt,bottom=6pt,before skip=6pt,after skip=6pt,#1}
\newtcolorbox{solutionbox}[1][]{enhanced,breakable,sharp corners,boxrule=0pt,frame hidden,colback=solvebg,left=8pt,right=8pt,top=6pt,bottom=6pt,before skip=6pt,after skip=6pt,#1}
\newtcolorbox{warningbox}[1][]{enhanced,breakable,sharp corners,boxrule=0pt,frame hidden,colback=commentbg,left=8pt,right=8pt,top=6pt,bottom=6pt,before skip=6pt,after skip=6pt,#1}
\hypersetup{pdfauthor={ИТМО},pdftitle={Лабораторная работа № 3-4. Две стадии}}
